\section{Mô hình nền tảng DNABERT-2}
\label{sec:dnabert2_detail}

DNABERT-2 \cite{zhou2023dnabert} được sử dụng làm backbone của PhaBERT-CNN. Phần này tóm tắt các đặc điểm kỹ thuật quan trọng của DNABERT-2 liên quan đến ứng dụng phân loại lối sống phage.

\subsection{Tokenization với BPE}

DNABERT-2 sử dụng Byte Pair Encoding (BPE) \cite{sennrich2015neural} thông qua SentencePiece \cite{kudo2018sentencepiece}:

\begin{itemize}
    \item \textbf{Vocabulary size}: 4,096 tokens
    \item \textbf{Average token length}: ~5 nucleotides
    \item \textbf{Sequence length reduction}: ~5× so với chuỗi gốc
    \item \textbf{Special tokens}: [CLS], [SEP], [PAD], [MASK], [UNK]
\end{itemize}

\textbf{Ví dụ tokenization:}
\begin{verbatim}
Input:  ATCGATCGATCGATCG (16 bp)
Tokens: [CLS] ATCG ATCG ATCG ATCG [SEP] (6 tokens)
\end{verbatim}

\subsection{Model architecture}

\begin{itemize}
    \item \textbf{Layers}: 12 transformer encoder layers
    \item \textbf{Hidden size}: 768 dimensions
    \item \textbf{Attention heads}: 12 heads (64 dimensions per head)
    \item \textbf{Intermediate size}: 3,072 (FFN hidden layer)
    \item \textbf{Total parameters}: ~117M
    \item \textbf{Activation}: GELU
    \item \textbf{Normalization}: Layer Normalization
    \item \textbf{Dropout}: 0.1
\end{itemize}

\subsection{Positional encoding: ALiBi}

DNABERT-2 sử dụng Attention with Linear Biases (ALiBi) \cite{press2021train} thay vì learned positional embeddings:

\begin{equation}
\text{Attention}(\mathbf{Q}, \mathbf{K}, \mathbf{V}) = \text{softmax}\left(\frac{\mathbf{Q}\mathbf{K}^T}{\sqrt{d_k}} + \mathbf{B}\right)\mathbf{V}
\end{equation}

với bias matrix:
\begin{equation}
\mathbf{B}_{ij} = m \cdot (j - i)
\end{equation}

trong đó $m$ là slope khác nhau cho mỗi attention head.

\textbf{Lợi ích cho phân loại phage:}
\begin{itemize}
    \item Xử lý contigs có độ dài biến thiên (100-1800 bp) mà không cần retrain
    \item Extrapolate tốt sang chuỗi dài hơn hoặc ngắn hơn training sequences
\end{itemize}

\subsection{Pre-training}

\textbf{Corpus:}
\begin{itemize}
    \item Multi-species genome corpus (bacteria, archaea, viruses, eukaryotes)
    \item Average sequence length: ~700 bp
    \item Masked Language Modeling objective
\end{itemize}

\textbf{Masking strategy:}
\begin{itemize}
    \item 15\% tokens được mask
    \item 80\%: Thay bằng [MASK]
    \item 10\%: Thay bằng random token
    \item 10\%: Giữ nguyên
\end{itemize}

\subsection{Output representations}

DNABERT-2 tạo ra ba loại representations:

\begin{enumerate}
    \item \textbf{Token-level embeddings}: $\mathbf{H} \in \mathbb{R}^{L \times 768}$ - Contextualized embedding cho mỗi token
    \item \textbf{[CLS] embedding}: $\mathbf{h}_{\text{CLS}} \in \mathbb{R}^{768}$ - Sequence-level representation
    \item \textbf{Pooled output}: Linear transformation của [CLS] embedding
\end{enumerate}

PhaBERT-CNN sử dụng \textbf{token-level embeddings} $\mathbf{H}$ làm input cho các lớp downstream, cho phép trích xuất đặc trưng chi tiết hơn so với chỉ sử dụng [CLS] embedding.
